\documentclass[11pt]{article}

\usepackage[utf8]{inputenc}
\usepackage[T1]{fontenc}
\usepackage[margin=1in]{geometry}
\usepackage{amsmath, amssymb, amsthm}
\usepackage{graphicx}
\usepackage{hyperref}
\usepackage{url}
\usepackage{booktabs}
\usepackage{xcolor}
\usepackage{microtype}
\usepackage{subcaption}
\usepackage{authblk}

\title{Spanda: Zero-Cost Lexical Entropy Matches Neural Semantic Uncertainty\\---Until Frontier Models Break It}

\author{
  Bhupen Nayak \\
  Independent Researcher \\
  \texttt{bhupennayak@icloud.com} \\
}

\begin{document}

\maketitle

\begin{abstract}
Detecting hallucinations in Large Language Models (LLMs) is critical for safe deployment, yet the dominant method---Semantic Entropy (SE)---requires expensive $O(K^2)$ neural NLI cross-encoding per query. We evaluate \textbf{Spanda} ($R_{sc}$), a zero-parameter, zero-latency exact-match normalized entropy metric, across four model scales: Qwen-1.5B, Mistral-7B, Qwen-3.6-27B, and GPT-OSS-120B, on GSM8K and TriviaQA.

We report two novel empirical findings. \textbf{(1) The Coherence Scaling Law:} As model capacity increases from 1.5B to 27B parameters, the AUROC of $R_{sc}$ for error detection rises monotonically from 0.577 to \textbf{0.889} ($p = 1.89 \times 10^{-28}$), matching and then surpassing heavy DeBERTa-v3 NLI Semantic Entropy---rendering neural cross-encoders redundant at the 7B--27B scale. \textbf{(2) Confident Mode Collapse:} On the 120B frontier model, we discover a critical failure mode where the model hallucinates the \textit{exact same incorrect answer} across all $K$ sampling paths, producing an \textit{inverted} AUROC of \textbf{0.091} (Cohen's $d = -2.23$, $p = 8.28 \times 10^{-15}$) on factual QA---demonstrating that self-consistency methods are \textit{fundamentally dangerous} on ungrounded frontier models. Together, these findings define the operational envelope for zero-cost uncertainty quantification: a $\sim$90,000$\times$ speedup over neural SE within a well-characterized safety boundary.
\end{abstract}

%----------------------------------------------------------------------
\section{Introduction}
%----------------------------------------------------------------------

As Large Language Models (LLMs) are deployed in safety-critical systems---from medical diagnosis to autonomous code generation---reliable detection of \textit{epistemic uncertainty} (``does the model know what it doesn't know?'') is essential. The dominant framework for epistemic uncertainty quantification is \textbf{Semantic Entropy (SE)} \cite{kuhn2023semantic}, which samples $K$ generation paths from the model, clusters semantically equivalent outputs using a secondary NLI cross-encoder (typically DeBERTa-v3-base), and measures the entropy over the resulting semantic clusters. When the model is uncertain, it generates diverse semantic meanings, yielding high entropy.

While principled, SE introduces two critical bottlenecks:
\begin{enumerate}
    \item \textbf{Quadratic NLI cost:} Clustering $K$ samples requires $\binom{K}{2}$ pairwise neural entailment classifications. For $K=10$, this is 45 GPU forward passes per query.
    \item \textbf{Latency:} Each DeBERTa-v3 inference adds $\sim$2 ms on a modern GPU. The total per-query overhead is $\sim$92 ms---prohibitive for real-time serving at scale.
\end{enumerate}

A natural question arises: \textit{is the neural NLI step always necessary?} Self-Consistency \cite{wang2022self} demonstrated that majority-vote over sampled chain-of-thought paths improves accuracy. But the field has assumed that lexical diversity is a poor proxy for semantic diversity, because LLMs can express the same meaning in many different surface forms.

\textbf{Our contribution.} We challenge this assumption empirically. We evaluate \textbf{Spanda} ($R_{sc}$), a zero-parameter exact-match normalized entropy metric, across four model scales spanning two orders of magnitude (1.5B to 120B parameters) on both mathematical reasoning (GSM8K) and factual recall (TriviaQA). We discover:

\begin{itemize}
    \item \textbf{The Coherence Scaling Law:} Larger models naturally converge on consistent lexical representations for correct answers, causing exact-match entropy to match neural SE at $\sim$7B and surpass it at 27B. This renders the expensive NLI step redundant for structured reasoning tasks on modern models.
    \item \textbf{Confident Mode Collapse:} At the extreme frontier (120B), we observe a previously undocumented failure mode: the model's parametric memory is so strong that it confidently hallucinates the \textit{same incorrect answer} across all $K$ paths, producing zero entropy and \textit{inverted} uncertainty scores. This finding has direct implications for RLHF alignment and self-consistency-based verification pipelines.
\end{itemize}

%----------------------------------------------------------------------
\section{Related Work}
%----------------------------------------------------------------------

\paragraph{Self-Consistency.} Wang et al.~\cite{wang2022self} showed that sampling multiple chain-of-thought reasoning paths and taking a majority vote improves accuracy over greedy decoding. However, they did not study whether the \textit{disagreement} among sampled paths can serve as a calibrated uncertainty signal.

\paragraph{Semantic Entropy.} Kuhn et al.~\cite{kuhn2023semantic} formalized epistemic uncertainty estimation by clustering sampled outputs using bidirectional NLI entailment. Their method achieves state-of-the-art AUROC for separating correct from incorrect predictions, but at significant computational cost. Farquhar et al.~\cite{farquhar2024detecting} extended this to conformal prediction sets. Neither work studied the interaction between model scale and the necessity of neural clustering.

\paragraph{Verbalized Confidence.} Kadavath et al.~\cite{kadavath2022language} explored prompting models to express their own confidence. However, verbalized confidence is poorly calibrated and easily gamed by instruction tuning.

\paragraph{Token-Level Entropy.} Lin et al.~\cite{lin2023generating} used per-token log-probabilities for uncertainty. This requires logit access (unavailable for most API-served models) and conflates lexical with semantic uncertainty.

\textbf{Gap.} No prior work has systematically studied how model scale affects the gap between cheap lexical clustering and expensive neural semantic clustering for uncertainty estimation. Our work fills this gap with a controlled four-tier evaluation.

%----------------------------------------------------------------------
\section{Methodology}
%----------------------------------------------------------------------

\subsection{The Spanda Metric ($R_{sc}$)}

Let $K$ be the number of sampled generations for a prompt $x$ at temperature $T > 0$. We extract final answers $y_1, \dots, y_K$ and partition them into $n$ equivalence classes $\{C_1, \dots, C_n\}$ via strict lexical matching after deterministic normalization (lowercasing, stripping whitespace/punctuation, standardizing numeric formatting).

The empirical probability of the $i$-th cluster is:
\begin{equation}
    w_i = \frac{|C_i|}{K}, \quad \sum_{i=1}^{n} w_i = 1
\end{equation}

The normalized Shannon entropy over this discrete distribution is:
\begin{equation}
    H_{\text{norm}} = \begin{cases}
        0 & \text{if } n = 1 \\
        \displaystyle\frac{-\sum_{i=1}^{n} w_i \ln w_i}{\ln K} & \text{if } n > 1
    \end{cases}
\end{equation}

The Spanda risk score combines entropy dispersion with modal dominance:
\begin{equation}
    R_{sc} = \alpha \cdot H_{\text{norm}} + (1 - \alpha) \cdot (1 - w_{\max})
\end{equation}
where $w_{\max} = \max_i w_i$ is the weight of the dominant cluster and $\alpha = 0.5$. $R_{sc} = 0$ indicates perfect consensus (high confidence); $R_{sc} \to 1$ indicates total disagreement (high uncertainty).

\subsection{NLI Semantic Entropy Baseline}

For direct comparison, we implement full Semantic Entropy following \cite{kuhn2023semantic}. For each question, $K$ sampled outputs are clustered via pairwise bidirectional entailment using \texttt{cross-encoder/nli-deberta-v3-base}. Two outputs are grouped if either entails the other. Normalized entropy is then computed over the resulting semantic clusters. This serves as the gold-standard neural baseline.

\subsection{Experimental Setup}

\begin{itemize}
    \item \textbf{Benchmarks:} GSM8K \cite{cobbe2021training} (mathematical reasoning, $\text{MAX\_NEW\_TOKENS} = 512$) and TriviaQA \cite{joshi2017triviaqa} (factual recall, $\text{MAX\_NEW\_TOKENS} = 64$).
    \item \textbf{Models:} Qwen-2.5-1.5B-Instruct (1.5B), Mistral-7B-Instruct-v0.1 (7B), Qwen-3.6-27B (27B), GPT-OSS-120B (120B).
    \item \textbf{Sampling:} $K = 10$ paths at $T = 0.7$ for 1.5B/7B (local GPU); $K = 5$ at $T = 0.7$ for 27B/120B (API-served via Groq).
    \item \textbf{Statistics:} Welch's $t$-test, AUROC (non-parametric concordance), Cohen's $d$, Bonferroni correction.
\end{itemize}

%----------------------------------------------------------------------
\section{Results}
%----------------------------------------------------------------------

\subsection{Multi-Scale Benchmark Results}

\begin{table*}[t]
\centering
\caption{Spanda ($R_{sc}$) vs.\ NLI Semantic Entropy (SE) across model scales. $^*$Rows marked with asterisk exhibit Confident Mode Collapse (see \S\ref{sec:mode_collapse}). All $p$-values: two-sided Welch's $t$-test.}
\label{tab:main_results}
\vspace{0.5em}
\resizebox{\textwidth}{!}{%
\begin{tabular}{llccccccc}
\toprule
\textbf{Model} & \textbf{Benchmark} & \textbf{$N$} & \textbf{Accuracy} & \textbf{$R_{sc}$ AUROC} & \textbf{SE AUROC} & \textbf{Cohen's $d$} & \textbf{$\Delta\bar{R}_{sc}$} & \textbf{$p$-value} \\
\midrule
\multicolumn{9}{l}{\textit{Small Scale (1.5B parameters)}} \\
Qwen-1.5B & GSM8K & 1{,}319 & 11.4\% & 0.577 & 0.584 & 0.27 & 0.072 & $5.0 \times 10^{-2}$ \\
Qwen-1.5B & TriviaQA & 500 & 32.0\% & 0.797 & \textbf{0.801} & 1.25 & 0.367 & $3.6 \times 10^{-25}$ \\
\midrule
\multicolumn{9}{l}{\textit{Mid Scale (7B parameters)}} \\
Mistral-7B & GSM8K & 1{,}319 & 8.2\% & \textbf{0.706} & 0.705 & 0.88 & 0.180 & $1.6 \times 10^{-7}$ \\
Mistral-7B & TriviaQA & 500 & 45.0\% & 0.698 & \textbf{0.755} & 0.70 & 0.216 & $1.1 \times 10^{-12}$ \\
\midrule
\multicolumn{9}{l}{\textit{Frontier Scale (27B--120B parameters)}} \\
Qwen-27B & GSM8K & 250 & 61.2\% & \textbf{0.889} & --- & \textbf{1.43} & 0.346 & $\mathbf{1.89 \times 10^{-28}}$ \\
Qwen-27B$^*$ & TriviaQA & 250 & 4.8\% & 0.344$^*$ & --- & $-$0.62 & $-$0.219 & $4.69 \times 10^{-2}$ \\
GPT-120B$^*$ & GSM8K & 250 & 92.8\% & 0.696$^*$ & --- & 1.58 & 0.251 & $5.40 \times 10^{-3}$ \\
GPT-120B$^*$ & TriviaQA & 250 & 10.4\% & 0.091$^*$ & --- & $-$2.23 & $-$0.467 & $8.28 \times 10^{-15}$ \\
\bottomrule
\end{tabular}%
}
\end{table*}

Table~\ref{tab:main_results} presents our main results. $\Delta\bar{R}_{sc}$ denotes the mean $R_{sc}$ separation between incorrect and correct predictions (positive = correct answers have lower $R_{sc}$, as expected). Several key patterns emerge.

\subsection{Finding 1: The Coherence Scaling Law}

On GSM8K (mathematical reasoning), $R_{sc}$ AUROC scales monotonically with model capacity:
\begin{equation}
    \text{AUROC}_{\text{GSM8K}}: \underbrace{0.577}_{\text{1.5B}} \longrightarrow \underbrace{0.706}_{\text{7B}} \longrightarrow \underbrace{0.889}_{\text{27B}}
    \label{eq:scaling}
\end{equation}

\textbf{Interpretation.} Small models (1.5B) produce syntactically noisy outputs---the same correct numerical answer may be formatted as ``42'', ``42.0'', ``the answer is 42'', or ``\$42''---breaking exact-match clustering and degrading $R_{sc}$. As model capacity increases, internal reasoning coherence improves: when a 27B model knows the answer, it consistently produces the same lexical token sequence across all $K$ paths. When it does not know, it produces genuinely diverse incorrect answers. This natural syntactic convergence means the expensive NLI step adds no discriminative value.

\textbf{Crossover point.} At Mistral-7B, $R_{sc}$ achieves an AUROC of \textbf{0.706}, matching and marginally surpassing the full NLI Semantic Entropy baseline (\textbf{0.705}). This crossover occurs because 7B models are syntactically coherent enough for mathematical answers but still produce occasional paraphrases for factual QA (where SE retains its advantage: 0.755 vs.\ 0.698).

\subsection{Finding 2: Confident Mode Collapse}
\label{sec:mode_collapse}

Our most striking finding emerges at the 120B frontier scale. On TriviaQA, GPT-OSS-120B achieves an \textit{inverted} AUROC of \textbf{0.091}---meaning $R_{sc}$ is \textit{anti-correlated} with correctness. The effect size is massive: Cohen's $d = -2.23$, $p = 8.28 \times 10^{-15}$.

\textbf{Mechanism.} When the 120B model encounters a factual question for which its parametric memory is insufficient, it does not exhibit the expected behavior of generating diverse guesses (which would yield high $R_{sc}$, correctly flagging uncertainty). Instead, its immense scale and RLHF-tuned confidence cause it to \textit{collapse} onto a single, highly confident hallucination that it reproduces identically across all $K$ sampling paths. This yields $R_{sc} = 0$ (perfect consensus), indistinguishable from genuinely correct answers.

Conversely, on the minority of questions the 120B model answers correctly (10.4\%), its correct answers occasionally exhibit minor formatting variation across paths, yielding \textit{higher} $R_{sc}$ than the confidently-hallucinated wrong answers.

\textbf{Quantitative evidence:}
\begin{itemize}
    \item Correct answers: $\bar{R}_{sc} = 0.564$ (moderate variation---the model hedges even when right)
    \item Incorrect answers: $\bar{R}_{sc} = 0.097$ (near-zero---the model is maximally ``confident'' in its hallucination)
\end{itemize}

This pattern is absent at the 27B scale, where incorrect answers correctly produce $\bar{R}_{sc} = 0.794$ (high uncertainty) and correct answers produce $\bar{R}_{sc} = 0.434$ (low uncertainty).

\textbf{Implication for AI safety.} This finding constitutes a direct challenge to self-consistency-based verification pipelines \cite{wang2022self}. Any system that uses ``do the model's multiple outputs agree?'' as a proxy for correctness will be \textit{systematically fooled} by frontier models on ungrounded factual queries. The failure is not stochastic---it is \textit{structurally guaranteed} by the model's parametric over-confidence.

\subsection{Finding 3: Latency Advantage}

For $K=10$ samples, NLI SE requires $\binom{10}{2} = 45$ DeBERTa-v3 forward passes at $\sim$2.05 ms each, totaling $\sim$92.4 ms per query. $R_{sc}$ reduces to $K$ string comparisons ($< 1\ \mu$s total in optimized implementations), yielding a $\sim$\textbf{90,000$\times$} latency reduction. At 10,000 queries per second, this saves $\sim$924 GPU-seconds per second of wall time---a massive cost reduction for production serving.

%----------------------------------------------------------------------
\section{Discussion}
%----------------------------------------------------------------------

\subsection{When to Use Spanda}

Our results delineate a clear operational envelope:

\begin{itemize}
    \item \textbf{Use $R_{sc}$} for structured reasoning tasks (math, code, extractive QA) on models in the \textbf{7B--27B range}. In this regime, $R_{sc}$ matches or exceeds neural SE at zero cost.
    \item \textbf{Do not use $R_{sc}$ alone} for open-domain factual QA on frontier ($>$100B) models. Confident Mode Collapse makes self-consistency signals \textit{anti-correlated} with correctness. External grounding (RAG) or hybrid approaches are required.
    \item \textbf{Use neural SE} when deploying small models ($<$7B) on tasks requiring semantic paraphrase detection (e.g., free-form QA where correct answers have multiple valid surface forms).
\end{itemize}

\subsection{Implications for RLHF and Alignment}

Confident Mode Collapse is likely an artifact of RLHF fine-tuning, which rewards models for sounding confident and helpful. This training signal inadvertently penalizes the natural epistemic hedging that smaller, less-tuned models exhibit. Our finding suggests that alignment procedures should explicitly preserve or incentivize output diversity on uncertain queries---a form of ``calibrated humility''---rather than optimizing solely for confident, helpful-sounding responses.

\subsection{Limitations}

\begin{itemize}
    \item \textbf{API-served models:} Our 27B and 120B evaluations use $K=5$ (vs.\ $K=10$ for locally-run models) due to API cost constraints. While $N > 240$ provides strong statistical power, a larger $K$ might modulate the Mode Collapse effect.
    \item \textbf{Model coverage:} We evaluate four model families. Extending to Claude, GPT-4, and Llama-3 would strengthen generalization claims.
    \item \textbf{Free-form text:} $R_{sc}$ is inherently limited to tasks with extractable, discrete answers. For open-ended generation, lightweight learned similarity functions remain necessary.
\end{itemize}

%----------------------------------------------------------------------
\section{Conclusion}
%----------------------------------------------------------------------

We have presented a systematic multi-scale evaluation of Spanda ($R_{sc}$), a zero-cost exact-match entropy metric for epistemic uncertainty quantification in LLMs. Our work reveals a \textbf{two-phase scaling law}: as model capacity grows from 1.5B to 27B, generation syntax converges sufficiently that cheap lexical clustering replaces expensive neural NLI---but at the extreme frontier (120B), structural over-confidence induces \textit{Confident Mode Collapse}, where hallucinated answers achieve perfect consensus and bypass all self-consistency checks. These findings jointly define when zero-cost uncertainty estimation is safe to deploy and when it is \textit{dangerous}, providing practitioners with a principled framework for high-throughput hallucination filtering in production LLM pipelines.

\paragraph{Reproducibility.} Code, evaluation scripts, and all raw experimental traces are publicly available at \url{https://github.com/Zach-al/Spnda}. The Spanda metric is released as a Python package installable via \texttt{pip install spanda}.

%----------------------------------------------------------------------
% References
%----------------------------------------------------------------------
\begin{thebibliography}{10}

\bibitem{kuhn2023semantic}
Kuhn, L., Gal, Y., and Farquhar, S.
\newblock Semantic uncertainty: Linguistic invariances for uncertainty estimation in natural language generation.
\newblock In \textit{International Conference on Learning Representations (ICLR)}, 2023.

\bibitem{wang2022self}
Wang, X., Wei, J., Schuurmans, D., Le, Q., Chi, E., Narang, S., Chowdhery, A., and Zhou, D.
\newblock Self-consistency improves chain of thought reasoning in language models.
\newblock In \textit{International Conference on Learning Representations (ICLR)}, 2023.

\bibitem{farquhar2024detecting}
Farquhar, S., Kossen, J., Kuhn, L., and Gal, Y.
\newblock Detecting hallucinations in large language models using semantic entropy.
\newblock \textit{Nature}, 630:625--630, 2024.

\bibitem{kadavath2022language}
Kadavath, S., et al.
\newblock Language models (mostly) know what they know.
\newblock \textit{arXiv preprint arXiv:2207.05221}, 2022.

\bibitem{lin2023generating}
Lin, S., Hilton, J., and Evans, O.
\newblock Teaching models to express their uncertainty in words.
\newblock \textit{Transactions on Machine Learning Research (TMLR)}, 2022.

\bibitem{cobbe2021training}
Cobbe, K., Kosaraju, V., Bavarian, M., Chen, M., Jun, H., Kaiser, L., Plappert, M., Tworek, J., Hilton, J., Nakano, R., Hesse, C., and Schulman, J.
\newblock Training verifiers to solve math word problems.
\newblock \textit{arXiv preprint arXiv:2110.14168}, 2021.

\bibitem{joshi2017triviaqa}
Joshi, M., Choi, E., Weld, D., and Zettlemoyer, L.
\newblock TriviaQA: A large scale distantly supervised challenge dataset for reading comprehension.
\newblock In \textit{Proceedings of ACL}, 2017.

\bibitem{belnap1977useful}
Belnap, N.~D.
\newblock A useful four-valued logic.
\newblock In \textit{Modern Uses of Multiple-Valued Logic}, Springer, 1977.

\bibitem{ouyang2022training}
Ouyang, L., et al.
\newblock Training language models to follow instructions with human feedback.
\newblock In \textit{Advances in Neural Information Processing Systems (NeurIPS)}, 2022.

\bibitem{touvron2023llama}
Touvron, H., et al.
\newblock Llama 2: Open foundation and fine-tuned chat models.
\newblock \textit{arXiv preprint arXiv:2307.09288}, 2023.

\end{thebibliography}

\end{document}
